\documentclass{amsart}
% For dealing with references we use the comment environment
\usepackage{verbatim}
\newenvironment{reference}{\comment}{\endcomment}
%\newenvironment{reference}{}{}
\newenvironment{slogan}{\comment}{\endcomment}
\newenvironment{history}{\comment}{\endcomment}

% For commutative diagrams we use Xy-pic
\usepackage[all]{xy}

% We use 2cell for 2-commutative diagrams.
\xyoption{2cell}
\UseAllTwocells

% We use multicol for the list of chapters between chapters
\usepackage{multicol}

% This is generally recommended for better output
\usepackage{lmodern}
\usepackage[T1]{fontenc}

% For cross-file-references

% Package for hypertext links:
\usepackage{hyperref}

% For any local file, say "hello.tex" you want to link to please

% Theorem environments.
%
\theoremstyle{plain}
\newtheorem{theorem}[subsection]{Theorem}
\newtheorem{proposition}[subsection]{Proposition}
\newtheorem{lemma}[subsection]{Lemma}

\theoremstyle{definition}
\newtheorem{definition}[subsection]{Definition}
\newtheorem{example}[subsection]{Example}
\newtheorem{exercise}[subsection]{Exercise}
\newtheorem{situation}[subsection]{Situation}

\theoremstyle{remark}
\newtheorem{remark}[subsection]{Remark}
\newtheorem{remarks}[subsection]{Remarks}

\numberwithin{equation}{subsection}

% Macros
%
\def\lim{\mathop{\mathrm{lim}}\nolimits}
\def\colim{\mathop{\mathrm{colim}}\nolimits}
\def\Spec{\mathop{\mathrm{Spec}}}
\def\Hom{\mathop{\mathrm{Hom}}\nolimits}
\def\Ext{\mathop{\mathrm{Ext}}\nolimits}
\def\SheafHom{\mathop{\mathcal{H}\!\mathit{om}}\nolimits}
\def\SheafExt{\mathop{\mathcal{E}\!\mathit{xt}}\nolimits}
\def\Sch{\mathit{Sch}}
\def\Mor{\mathop{\mathrm{Mor}}\nolimits}
\def\Ob{\mathop{\mathrm{Ob}}\nolimits}
\def\Sh{\mathop{\mathit{Sh}}\nolimits}
\def\NL{\mathop{N\!L}\nolimits}
\def\CH{\mathop{\mathrm{CH}}\nolimits}
\def\proetale{{pro\text{-}\acute{e}tale}}
\def\etale{{\acute{e}tale}}
\def\QCoh{\mathit{QCoh}}
\def\Ker{\mathop{\mathrm{Ker}}}
\def\Im{\mathop{\mathrm{Im}}}
\def\Coker{\mathop{\mathrm{Coker}}}
\def\Coim{\mathop{\mathrm{Coim}}}

% Boxtimes
%
\DeclareMathSymbol{\boxtimes}{\mathbin}{AMSa}{"02}

%
% Macros for moduli stacks/spaces
%
\def\QCohstack{\mathcal{QC}\!\mathit{oh}}
\def\Cohstack{\mathcal{C}\!\mathit{oh}}
\def\Spacesstack{\mathcal{S}\!\mathit{paces}}
\def\Quotfunctor{\mathrm{Quot}}
\def\Hilbfunctor{\mathrm{Hilb}}
\def\Curvesstack{\mathcal{C}\!\mathit{urves}}
\def\Polarizedstack{\mathcal{P}\!\mathit{olarized}}
\def\Complexesstack{\mathcal{C}\!\mathit{omplexes}}
% \Pic is the operator that assigns to X its picard group, usage \Pic(X)
% \Picardstack_{X/B} denotes the Picard stack of X over B
% \Picardfunctor_{X/B} denotes the Picard functor of X over B
\def\Pic{\mathop{\mathrm{Pic}}\nolimits}
\def\Picardstack{\mathcal{P}\!\mathit{ic}}
\def\Picardfunctor{\mathrm{Pic}}
\def\Deformationcategory{\mathcal{D}\!\mathit{ef}}

\setlength{\emergencystretch}{3em}
\begin{document}
\title[Stacks Project, Tag 0BJ0]{350 --- Nagata local rings and geometrically reduced formal fibres}
\maketitle
\noindent Copyright (C) 2005--2025 Johan de Jong. Stacks Project authors. Source: \href{https://stacks.math.columbia.edu/tag/0BJ0}{Tag 0BJ0}.

\noindent Editorial difficulty note (not author text). Difficulty H2: One direction uses finite local domain extensions and analytic unramifiedness to control formal fibres. The reverse direction constructs a finite integral model inside an arbitrary finite residue-field extension and identifies its completed localization, so all geometric field extensions are addressed..

\noindent Editorial provenance note (not author text). History: excerpt from revision \texttt{a0444\allowbreak 6e57e\allowbreak c1fbc\allowbreak 252a8\allowbreak 71afc\allowbreak ec775\allowbreak 2fb28\allowbreak 07b14}, more-algebra.tex, lines 13412--13461. Reference presentation was adapted; original-author context and core support blocks were retained or added. The mathematical wording is unchanged. Licensed under GNU FDL 1.2 or later; no invariant sections or cover texts. See ../licenses/COPYING. Context blocks reproduce author definitions and supporting results; they are not additional counted problems.

\begin{lemma}
\label{lemma-Nagata-local-ring}
Let $(A, \mathfrak m)$ be a Noetherian local ring.
The following are equivalent
\begin{enumerate}
\item $A$ is Nagata, and
\item the formal fibres of $A$ are geometrically reduced.
\end{enumerate}
\end{lemma}

\begin{proof}
Assume (2). By
Algebra, Lemma \href{https://stacks.math.columbia.edu/tag/0BI2}{lemma local nagata and analytically unramified (Tag 0BI2)}
we have to show that if $A \to B$ is finite, $B$ is a domain,
and $\mathfrak m' \subset B$ is a maximal ideal, then $B_{\mathfrak m'}$
is analytically unramified.
Combining Lemmas \href{https://stacks.math.columbia.edu/tag/0BIW}{lemma formal fibres reduced (Tag 0BIW)} and
\href{https://stacks.math.columbia.edu/tag/0BIU}{lemma check P ring maximal ideals (Tag 0BIU)} and
Proposition \href{https://stacks.math.columbia.edu/tag/0BIV}{proposition finite type over P ring (Tag 0BIV)}
we see that the formal fibres of $B_{\mathfrak m'}$ are
geometrically reduced. In particular
$B_{\mathfrak m'}^\wedge \otimes_B L$ is reduced
where $L$ is the fraction field of $B$.
It follows that $B_{\mathfrak m'}^\wedge$ is reduced, i.e.,
$B_{\mathfrak m'}$ is analytically unramified.

\medskip\noindent
Assume (1). Let $\mathfrak q \subset A$ be a prime ideal
and let $K/\kappa(\mathfrak q)$ be a finite extension.
We have to show that $A^\wedge \otimes_A K$ is reduced.
Let $A/\mathfrak q \subset B \subset K$ be a local subring
finite over $A$ whose fraction field is $K$.
To construct $B$ choose $x_1, \ldots, x_n \in K$
which generate $K$ over $\kappa(\mathfrak q)$
and which satisfy monic polynomials
$P_i(T) = T^{d_i} + a_{i, 1} T^{d_i - 1} + \ldots + a_{i, d_i} = 0$
with $a_{i, j} \in \mathfrak m$. Then let $B$ be the $A$-subalgebra
of $K$ generated by $x_1, \ldots, x_n$. (For more details see
the proof of Algebra, Lemma
\href{https://stacks.math.columbia.edu/tag/0BI2}{lemma local nagata and analytically unramified (Tag 0BI2)}.)
Then
$$
A^\wedge \otimes_A K =
(A^\wedge \otimes_A B)_\mathfrak q =
B^\wedge_\mathfrak q
$$
Since $B^\wedge$ is reduced by  Algebra, Lemma
\href{https://stacks.math.columbia.edu/tag/0BI2}{lemma local nagata and analytically unramified (Tag 0BI2)}
the proof is complete.
\end{proof}
\end{document}
